% ECONOMICS_PROBLEM: {"id": "H21-528", "title": "Optimal VAT/consumption-tax design", "jel_code": "H21", "source_id": "OP-0044"}
% !TeX program = xelatex
\documentclass[11pt,a4paper]{article}
\usepackage[margin=23mm]{geometry}
\usepackage{fontspec}
\setmainfont{Latin Modern Roman}
\usepackage{amsmath,amssymb,mathtools}
\usepackage{xcolor,enumitem,booktabs,longtable,array,xurl,hyperref}
\definecolor{muted}{HTML}{53636C}
\hypersetup{colorlinks=true,urlcolor=blue,linkcolor=blue}
\setlength{\parindent}{0pt}
\setlength{\parskip}{5pt}
\newcommand{\R}{\mathbb R}
\newcommand{\E}{\mathbb E}
\newcommand{\N}{\mathbb N}
\newcommand{\ind}{\mathbf 1}
\newcommand{\Var}{\operatorname{Var}}
\newcommand{\Cov}{\operatorname{Cov}}
\newcommand{\supp}{\operatorname{supp}}
\DeclareMathOperator*{\argmin}{arg\,min}
\DeclareMathOperator*{\argmax}{arg\,max}
\newcommand{\entrymeta}[3]{{\sffamily\small\color{muted}\textbf{\texttt{#1}}\quad|\quad #2\par #3\par}}
\newcommand{\Problemref}[1]{\texttt{#1}}
\newcommand{\JELref}[2]{\texttt{#2} (JEL \texttt{#1})}
\begin{document}
\section*{Optimal VAT/consumption-tax design}
\textbf{Problem ID:} \texttt{H21-528}\quad \textbf{JEL:} \texttt{H21}\par
% BEGIN ECONOMICS STATEMENT
\label{problem:OP-0044}
\label{atlas:prob:P4}
\noindent\textbf{Original atlas identifier:} P4; entry 44. \textbf{Field:} Public finance and taxation. \textbf{Original tractability label:} Medium.

\noindent\textbf{Classification:} Parameterized theoretical/empirical research agenda; not a certified unsolved theorem.

\subsection*{Definitions and assumptions}
Fix households $i$ with labor productivity, preferences over $K$ consumption goods, and welfare weight $w_i\geq0$. A policy $p=(T,\boldsymbol\tau,\boldsymbol e)$ specifies a labor-income tax/transfer schedule $T$, commodity tax rates $\boldsymbol\tau=(\tau_1,\ldots,\tau_K)$, exemptions, refund eligibility, and enforcement intensities $\boldsymbol e$. A production/invoice-network model $m$ determines producer prices, pass-through, household demand, input use, evasion, and the selected equilibrium. Let $EV_i(p;m)$ be monetary equivalent variation, including compliance burdens; $R(p;m)$ is revenue net of lawful refunds, and $C(p;m)$ is administrative resource cost. Specify whether preferences are weakly separable between consumption and labor and whether the consumption subutility is common across productivity types; do not assume these restrictions implicitly.

\subsection*{Formal research question}
\emph{Editorial formulation: the mathematical target below is not a theorem quoted from the references.}

Given feasible income-tax instruments, legal VAT constraints, enforcement capacity, and a required net-revenue target $\bar R$, determine
\[
 \sup_{p\in\mathcal P:\ R(p;m)-C(p;m)\geq\bar R}\mathbb E_m[w_iEV_i(p;m)].
\]
Characterize when uniform positive commodity taxation, differential rates, or exemptions are $\varepsilon$-optimal in a specified model class. Identify the distributional and evasion effects of moving an item from exemption to creditable taxation; exemption and zero-rating must be treated as distinct instruments because input-tax recovery differs. With scanner and invoice data, determine whether demand, production substitution, and concealment responses are separately identified under declared variation and exclusion restrictions. If administrative records omit informal transactions, give sensitivity bounds based on explicit missing-sales assumptions.

\subsection*{Interpretation and scope}
The Atkinson--Stiglitz benchmark does not license a blanket statement that actual VAT differentiation is always inefficient: its conclusion is conditional on preference and income-tax assumptions. Nor does observing low household expenditure on a taxed item establish low welfare cost, because income effects, quality choice, and producer prices can change. Rate design should account for invoice-chain breaks, business entry, and refund fraud as well as household incidence. Use resource costs once: private compliance costs already entering equivalent variation must not be subtracted a second time. Household surveys and scanner data may miss the poorest or informal purchases, so sample coverage must accompany any redistribution claim. The extension is a joint enforceability/redistribution comparison with explicitly specified instruments and population coverage.

\subsection*{Source audit and status}
[S1] is a solved direct/indirect-tax benchmark, [S2] studies VAT adoption, and [S3] provides information/enforcement evidence. The European VAT webpage in the atlas is institutional background, not the mathematical source. The expanded optimization problem is editorial and lacks a cited universal unresolved theorem.

\subsection*{References}
\begin{enumerate}[label={[S\arabic*]},leftmargin=*,itemsep=0.6em]
\item Anthony B. Atkinson and Joseph E. Stiglitz (1976). The Design of Tax Structure: Direct versus Indirect Taxation. Journal of Public Economics 6(1--2), 55--75. DOI: 10.1016/0047-2727(76)90041-4.
\url{https://www.sciencedirect.com/science/article/abs/pii/0047272776900414}
\newline\emph{Locator and support limits:} Article: direct/indirect-tax benchmark under preference and information restrictions; enforcement failures are additional assumptions.
\item Michael Keen and Ben Lockwood (2010). The Value Added Tax: Its Causes and Consequences. Journal of Development Economics 92(2), 138--151. DOI: 10.1016/j.jdeveco.2009.01.012.
\url{https://www.sciencedirect.com/science/article/abs/pii/S0304387809000133}
\newline\emph{Locator and support limits:} Abstract: VAT adoption and revenue consequences; not causal welfare comparisons of every rate schedule.
\item Dina Pomeranz (2015). No Taxation without Information: Deterrence and Self-Enforcement in the Value Added Tax. American Economic Review 105(8), 2539--2569. DOI: 10.1257/aer.20130393.
\url{https://www.aeaweb.org/articles?id=10.1257/aer.20130393}
\newline\emph{Locator and support limits:} Abstract and experiments: Chilean VAT information and supply-chain enforcement; no comprehensive privacy-cost estimate.
\end{enumerate}

\subsection*{Common conventions: Source mathematical conventions}

Unless an entry states otherwise, agent and firm sets are finite. State and action spaces are measurable, strategies and outcome maps are measurable, probability laws are specified, and expectations exist for the targets under discussion. Infinite-horizon discount factors lie in $(0,1)$ and discounted objectives are finite. Norms, tolerances and complexity input sizes must be specified for computational comparisons. Constants or primitives that appear locally are inputs, not empirical facts asserted by this catalogue.

\textbf{Identification.} A statistical model $m\in\mathcal M$ implies an observable law $P_m$ and target $\theta(m)$. At law $P$, its identified set is
\[
\Theta(P;\mathcal M)=\{\theta(m):m\in\mathcal M,\ P_m=P\}.
\]
Point identification holds when this set is a singleton, or equivalently when $P_{m_1}=P_{m_2}$ implies $\theta(m_1)=\theta(m_2)$ for all compatible models. Sharp bounds mean bounds attained, or approached when the boundary is not attained, by compatible models. Writing this set is a definition, not a solution: the substantive task is to establish informative restrictions and derive or estimate the set. An unrestricted-confounding model may give only trivial bounds.

\textbf{Inference.} Identification, estimability and valid uncertainty quantification are separate questions. Fix $\alpha\in(0,1)$. A confidence procedure $C_n$ over a declared law class $\mathcal P$ has asymptotic uniform coverage at level $1-\alpha$ when
\[
\liminf_{n\to\infty}\inf_{P\in\mathcal P}\Pr_P\{\theta(P)\in C_n\}\geq1-\alpha.
\]
For set-identified targets the coverage event must be defined separately, for example coverage of the full identified set. If an entry uses identification-indexed models instead of uniquely indexed laws, the same coverage convention is applied over those models. Pointwise limits do not establish uniform validity, and asymptotic validity does not guarantee finite-sample precision. Sample sizes, dependence structures and weak-identification sequences are part of each application.

\textbf{Counterfactuals.} $Y_i(a)$ is a potential outcome under a specified intervention, including its timing, implementation and versions. It is not $Y_i$ conditional on voluntarily choosing $a$. A full intervention vector permits interference; shorthand scalar treatment notation does not silently impose no interference. Assignment, selection, attrition, anticipation and measurement must be specified. A claim about an instrument's local effect is not automatically a population policy effect.

\textbf{Economic equilibrium.} A counterfactual model includes best responses, constraints, feasibility, market clearing, state transitions and expectations consistent with the stated information. When several equilibria exist, either a declared selector is an input or the target is set-valued. Comparisons across policy regimes must use the stated selection convention. Reduced-form relationships are not presumed invariant after an equilibrium-changing intervention.

\textbf{Optimization and welfare.} Optimization is over the stated feasible set. If attainment is not established, $\sup$ or $\inf$ and $\varepsilon$-optimal choices replace an asserted maximizer or minimizer. For example, with value $V=\sup_{a\in\mathcal A}W(a)<\infty$, an $\varepsilon$-optimal set is $\{a\in\mathcal A:W(a)\geq V-\varepsilon\}$ for $\varepsilon>0$. Benefits, costs, risk and health or environmental outcomes can be aggregated only using declared units and conversion weights. Interpersonal utility weights, the treatment of fiscal transfers and foreign profits, discounting, and the behavioral welfare criterion are explicit inputs. A comparative-static derivative additionally requires existence and appropriate differentiability of the selected counterfactual.

\textbf{Sources.} Local labels [S1], [S2], and so forth restart in every question. Locators identify the relevant abstract, model, section or result at the level actually checked; a bibliographic or abstract-level verification is not represented as inspection of an entire proof. Working papers are distinguished from later publications and changing versions. Source summaries are paraphrased. All URL checks and editorial formulations refer to the audit date stated above.
% END ECONOMICS STATEMENT
\end{document}
